\Qhead{3.7}{Тепловой шум. Тепловое излучение абсолютно чёрного тела}{Тепловой шум; классический и квантовый варианты формулы Найквиста. Тепловое излучение абсолютно черного тела.}
\Subfile{3.7.1}{sections/03_statistics/q07/s01}
\Subfile{3.7.2}{sections/03_statistics/q07/s02}

\recap{Найквист связывает шум пассивного импеданса с температурой и потерями. Квантовая формула различает симметризованные флуктуации и излучаемую мощность. Планковский спектр получается из плотности мод и бозевского заселения.}
\board{$4k_BTR$ и проверка $k_BT/C$; квантовый множитель; плотность мод и спектр Планка.}
\oral{Можно ли получить ненулевую тепловую мощность от резистора при нулевой температуре за счёт $hf/2$?}
\bait{Можно оставить независимость согласованной шумовой мощности от сопротивления: делитель напряжения компенсирует величину шумовой ЭДС.}
\sources{Биленко и соавт.~\cite{bilenko}, §4.2.1; Clerk и соавт.~\cite{clerk}, §II и приложение~C: различие квантовых спектров. Планковская плотность выведена подсчётом мод.}
